\par\noindent\textbf{(1)} 设介质位移为 \(u\)，局域谐振器位移 \(v\)，弹簧刚度每体积 \(k_r=m_r\omega_r^2\)。谐振器方程 \(m_r\ddot v=-k_r(v-u)\)，代入 \(e^{-i\omega t}\) 得 \(v=\omega_r^2u/(\omega_r^2-\omega^2)\)。
\par\noindent\textbf{(2)} 宏观动量含背景质量与谐振器惯性，得 \(\rho_{\rm eff}(\omega)=\rho_0+m_r\omega_r^2/(\omega_r^2-\omega^2)\)。当 \(\omega\to0\) 时 \(\rho_{\rm eff}\to\rho_0+m_r\)，而在 \(\omega>\omega_r\) 附近动态惯性可为负。
\par\noindent\textbf{(3)} 有效色散 \(k^2=\omega^2\rho_{\rm eff}/K\)。令 \(\rho_{\rm eff}=0\) 得 \(\omega_z^2=\omega_r^2(1+m_r/\rho_0)\)。对 \(\omega_r<\omega<\omega_z\)，\(\rho_{\rm eff}<0\)、\(k^2<0\)，波在均匀有效介质内倏逝，形成局域共振禁带。
\par\noindent\textbf{(4)} 理想无耗散 \(\rho_{\rm eff}\) 在 \(\omega=\omega_r\) 有极点，此处局域内部振幅很大，线性近似及均匀化可能失效。若加阻尼 \(2\gamma\dot v\)，分母改为 \(\omega_r^2-\omega^2-2i\gamma\omega\)，有效密度变复数并出现吸收，而非真实无限响应。
